\section{SDE de varianza explosiva: VE-SDE}

\subsection{De SMLD a VE-SDE}

SMLD (Score Matching con Dinámica de Langevin, también conocido como NCSN) define la distribución de transición hacia adelante como:

$$
q(\mathbf{x}_t \mid \mathbf{x}_{t-1}) = \mathcal{N}\!\left(\mathbf{x}_t;\, \mathbf{x}_{t-1},\, (\sigma_t^2 - \sigma_{t-1}^2)\mathbf{I}\right)
$$

La distribución de salto correspondiente (de $\mathbf{x}_0$ directamente a $\mathbf{x}_t$) es:

$$
q(\mathbf{x}_t \mid \mathbf{x}_0) = \mathcal{N}\!\left(\mathbf{x}_t;\, \mathbf{x}_0,\, \sigma_t^2 \mathbf{I}\right)
$$

\begin{importantbox}{La idea central de SMLD}
El proceso hacia adelante en SMLD consiste esencialmente en \textbf{incrementar gradualmente la varianza del ruido gaussiano}, sin desplazar la media de los puntos de datos. Es decir, cada punto de datos $\mathbf{x}_0$ se "difumina": a medida que avanza el tiempo, la distribución gaussiana centrada en $\mathbf{x}_0$ se vuelve progresivamente más amplia, hasta que toda la distribución de datos se aproxima a una distribución gaussiana con gran varianza. En el proceso inverso, la función score $\nabla_{\mathbf{x}} \log p_t(\mathbf{x})$ guía el retorno desde lo difuso hacia lo nítido.
\end{importantbox}

\subsection{Derivación de VE-SDE}

Se reescribe la cadena de Markov discreta en forma continua. Cuando el número de pasos de tiempo $N \to \infty$ y el intervalo temporal $\Delta t \to 0$, la probabilidad de transición se convierte en:

$$
\mathbf{x}_{t+\Delta t} = \mathbf{x}_t + \sqrt{\sigma^2(t+\Delta t) - \sigma^2(t)} \cdot \boldsymbol{\epsilon}, \quad \boldsymbol{\epsilon} \sim \mathcal{N}(0, \mathbf{I})
$$

Utilizando el desarrollo de Taylor $\sigma^2(t+\Delta t) \approx \sigma^2(t) + \frac{\mathrm{d}[\sigma^2(t)]}{\mathrm{d}t}\Delta t$ y tomando el límite, se obtiene:

$$
\mathrm{d}\mathbf{x} = \sqrt{\frac{\mathrm{d}[\sigma^2(t)]}{\mathrm{d}t}}\;\mathrm{d}\mathbf{w}
$$

Esta es la \textbf{SDE de varianza explosiva (VE-SDE)}. Los coeficientes correspondientes son:

\begin{itemize}
    \item \textbf{Coeficiente de arrastre}: $f(\mathbf{x}, t) = 0$ (no hay arrastre determinista)
    \item \textbf{Coeficiente de difusión}: $g(t) = \sqrt{\dfrac{\mathrm{d}[\sigma^2(t)]}{\mathrm{d}t}}$
\end{itemize}

\begin{knowledgebox}{¿Por qué se llama "varianza explosiva"}
En la VE-SDE, dado que el término de arrastre es cero, la media permanece constante ($\mathbb{E}[\mathbf{x}_t \mid \mathbf{x}_0] = \mathbf{x}_0$), mientras que la varianza $\sigma^2(t)$ aumenta monótonicamente con el tiempo y no tiene cota superior (si no se realiza un recorte). Por ello se denomina Variance Exploding: la varianza crece de manera "explosiva". La imagen física correspondiente es: cada punto de datos permanece "quieto en su sitio", pero la nube gaussiana de ruido que lo rodea se vuelve progresivamente más grande.
\end{knowledgebox}

\subsection{$\alpha_t$ y $\sigma_t$ en VE-SDE}

Bajo el marco unificado de Variational Diffusion Models, la VE-SDE corresponde a:

$$
\alpha_t = 1, \qquad \sigma_t = \sigma(t)
$$

Donde $\sigma(t)$ es la función de escala del ruido definida en SMLD. La relación señal-ruido (SNR) es:

$$
\mathrm{SNR}(t) = \frac{\alpha_t^2}{\sigma_t^2} = \frac{1}{\sigma^2(t)}
$$

Dado que $\sigma(t)$ aumenta monótonamente, la SNR disminuye monótonamente, cumpliendo las condiciones de Variational Diffusion Models.

\subsection{Resumen de la sección}

La VE-SDE es una extensión natural de SMLD en el dominio del tiempo continuo: sin arrastre, solo difusión, con una varianza que aumenta continuamente a lo largo del tiempo. Corresponde al proceso de aumento gradual de la varianza del ruido en las Annealed Langevin Dynamics.
